% Raw abstract body text (no environment wrapper). The enclosing container
% in main.tex (tcolorbox header panel) handles formatting.
Scientific publication compresses a branching, iterative research process into a linear narrative, discarding the majority of what was discovered along the way. This compilation imposes two structural costs: a \textbf{Storytelling Tax}, where failed experiments, rejected hypotheses, and the branching exploration process are discarded to fit a linear narrative; and an \textbf{Engineering Tax}, where the gap between reviewer-sufficient prose and agent-sufficient specification leaves critical implementation details unwritten. Tolerable for human readers, these costs become critical when AI agents must understand, reproduce, and extend published work. We introduce the \textbf{Agent-Native Research Artifact (\ara{})}, a protocol that replaces the narrative paper with an agent-executable research package structured around four layers: scientific logic, executable code with full specifications, an exploration graph that preserves the failures compilation discards, and evidence grounding every claim in raw outputs. Three mechanisms support the ecosystem: a \textbf{Live Research Manager} that captures decisions and dead ends during ordinary development; an \textbf{\ara{} Compiler} that translates legacy PDFs and repos into \ara{}s; and an \textbf{\ara{}-native review system} that automates objective checks (analogous to a grammar checker for prose) so human reviewers can focus on significance, novelty, and taste. On PaperBench and RE-Bench, \ara{} raises question-answering accuracy from 72.4\% to 93.7\% and reproduction success from 57.4\% to 64.4\%. On RE-Bench's five open-ended extension tasks, preserved failure traces in \ara{} accelerate progress, but can also constrain a capable agent from stepping outside the prior-run box depending on the agent's capabilities.
